\documentclass{ctexbook}
\usepackage{amsmath,amssymb,amsthm}
\usepackage[hidelinks]{hyperref}
\newcommand{\Spec}{\operatorname{Spec}}
\newcommand{\ednote}[1]{\footnote{英译编者注：#1}}
\newcommand{\tnote}[1]{\par\smallskip\noindent\textbf{本批译校注：}#1\par\smallskip}
\newcommand{\entry}[2]{\par\bigskip\noindent\textbf{#1\quad #2}\quad}
\title{代数几何基础第五卷预备笔记\\第7节独立续译与译校说明}
\author{据用户上传英译编辑稿译出}
\date{2026年11月14日}
\begin{document}
\maketitle
\frontmatter
\chapter*{范围与底本}
本批完整翻译用户上传的 \texttt{EGAVaiello.pdf} 第一篇《超平面截面与锥投影》第7节，覆盖（7.1）至（7.5），包括原有证明、编写说明、难辨字标记与英译编者脚注26--28。底本节标题起于第18页末，正文见第19--21页。本批不是整部第五卷的全译，也没有覆盖第8节。

底本题为 Alexandre Grothendieck's EGA V，署为 Piotr Blass 与 Joseph Blass 对预备笔记的翻译和编辑。本译文依据这份用户上传的英文稿，并非声称已经对照法文手稿逐字校勘。保留原作者未定稿性质，保留疑点，不把疑似错误静默改成正确命题。原稿、英译编者注、本批译校注分别标识。

本批作为第三版合订本的独立增补，不覆盖已有合订本或工程。第三版现存进度记录中，这一篇已译至第6节。为避免与原对话可能正在处理的第0章及第四卷主线发生并发覆盖，暂不自动合入合订本。符号沿用前文：$P^\vee$ 为超平面参数空间，$Y_\xi$ 为相应截面；$F,G_\xi,\phi_\xi$ 的具体构造承前文。

\mainmatter
\chapter*{§7\quad 构造指定类型的超平面截面与多截面的应用}
\addcontentsline{toc}{chapter}{§7 构造指定类型的超平面截面与多截面的应用}
\entry{7.1}{}
注意，如果 $S=\Spec(k)$，其中 $k$ 是无限域，则 $P^\vee$ 的每个非空开子集都含有一个 $k$-有理点。因此，采用（4.1）的记号，如果由可构造性质 $P$ 定义的集合 $E$ 含有泛点 $\eta$，那么它就含有一个 $k$-有理点，从而存在一个定义在 $k$ 上、具有性质 $P$ 的 $X$ 的超平面截面。

另一方面，现在让 $S$ 再次是任意的。立即可见，对每个 $s\in S$，以及纤维 $P_s^\vee$ 中每个在 $k(s)$ 上有理的点 $\xi$，都存在 $s$ 的某个开邻域 $U$ 上的 $P^\vee$ 的截面 $\xi$，且它经过 $\xi_0$。
\tnote{底本这一句先把纤维中的点写为 $\xi$，末尾却写为 $\xi_0$，而 $\xi$ 也被用于截面。这里保留其区别未定的原符号；按下文的用法，纤维点应记为 $\xi_0$，截面记为 $\xi$。这只是符号辨析，不是另加数学结论。}
如果 $E$ 仍像（4.1）那样由可构造性质 $P$ 定义，而且已知 $E$ 为开集，例如这是第5节的结果，那么当 $\xi_0\in E$ 时，至少在把 $U$ 充分缩小以后，截面 $\xi$ 就是 $U$ 上取值于 $E$ 的截面。因此，可以在 $s$ 的一个开邻域 $U$ 上构造 $X$ 的超平面截面 $Y_\xi$，使对每个 $t\in U$，它在 $t$ 处的纤维 $Y_{\xi(t)}$ 都满足性质 $P$。如果事先没有给定 $\xi_0$，但 $k(s)$ 是无限域，那么可以结合以上两点，在 $s$ 的一个开邻域上得到具有上述性质的超平面截面。

最后，利用第5节，可以得到一个判据，以断言 $Y_\xi$，分别地 $G_\xi$，也在 $S$ 上平坦。这里假设 $X$，分别地 $F$，在 $S$ 上平坦，以便应用该处的结果。因此，以 $Y_\xi$ 代替 $X$，就能迭代上述构造；例如，在某些条件下，可以在给定点 $s$ 的一个开邻域 $U$ 上逐步构造（英译稿写作“越来越接近，通过逐次逼近？？？”）\ednote{法文为 \textit{de proche en proche}。〔底本脚注26。〕}一个 $X$ 的“多截面”$S'$，使得
\[
 S'\longrightarrow U
\]
有限、平坦，而且其纤维满足性质 $P$。

如果 $k(s)$ 是有限域，则在能够应用前面的构造之前，也许必须先作一个平展且满射的基变换 $S'\to U$，其中 $U$ 是 $s$ 的开邻域。事实上，在（6.1）开头的条件下，如果 $k$ 有限，则非空开集 $U$ 中未必存在 $k$-有理点，但 $U$ 中必定有闭点，因而有取值于 $k$ 的某个有限扩张 $k'$ 的点；这个扩张必为可分扩张。当 $k=k(s)$ 时，因此可以在 $s$ 的某个邻域 $U$ 上作适当的有限平展扩张 $S'$，对应于剩余域扩张 $k'$，即使得
\[
 S_s^1\xrightarrow{\sim}\Spec(k'),
\]
然后经基变换 $S'\to S$，归结到 $S'$ 中位于 $s$ 上方的唯一点 $s'$ 所处的有利情形。
\tnote{上式底本确实排作 $S_s^1$，而上下文使用 $S'$；疑为 $S'_s$ 的排印问题，未静默替换。该段又回指（6.1）开头，而那里假设 $k$ 代数闭，此处却讨论 $k$ 有限；回指存在不一致，可能原意是指（7.1）开头，但未据猜测改写。$U$ 在段中同时用于基底邻域及参数开集，合并定稿时须区别。}

不过，对于（4.2）和（4.3），我还须指出一项遗憾〔法文：\textit{un remords}〕：它们本应以稍广一些的形式陈述，至少作为一条备注。给定整数 $m$，令 $E$ 为那些 $\xi\in P^\vee$ 所成的集合，使得 $G_\xi$，分别地 $Y_\xi$，除去一个维数 $\leq m$ 的闭集以外满足性质 $P$，即 $P$-奇异点集 $Z$ 的维数 $\leq m$。那么：
\begin{enumerate}
\item[(a)] $E$ 是 $P^\vee$ 的可构造子集；
\item[(b)] 当 $S=\Spec(k)$ 时，如果 $F$，分别地 $X$，除去一个维数 $\leq m+1$ 的集合以外满足 $P$，则 $E$ 包含一个非空开集。
\end{enumerate}

\entry{7.2}{命题}
假设 $X$ 在 $S$ 上固有，$F$ 为有限表示并在 $S$ 上平坦。\footnote{本批译校注：英文正文为 ``of finite presentation finite and flat over $S$''，其中第二个 ``finite'' 的语法关系不明。未据这个重复的词添加独立有限性假设。}
设 $P$ 为（4.2）中(i)至(v)之一，$m$ 为整数。取 $S\in S$，并假设 $X_s$ 中 $F_s$ 不满足 $P$ 的点的集合 $Z_s$ 的维数 $\leq m+1$。如果 $k(s)$ 还是无限域，则在 $E$ 中存在 $s$ 的一个邻域 $U$，以及 $U$ 上 $P^\vee$ 的一个截面 $\xi$，具有下列性质：对每个 $x\in U$，$Y_{\xi,s}$ 中 $G_{\xi(s)}$ 不满足 $P$ 的点的集合的维数 $\leq m$，而且 $\phi_{\xi(s)}$ 是 $F_{\xi,s}$-正则的。在这些条件下，$Y_\xi$ 上的模 $G_\xi$ 相对 $U$ 平坦。最后，即使不假设 $k(s)$ 无限，也仍可以在作（7.1）所预期类型的平展扩张后完成上述构造。
\tnote{本段保留了底本的 $S\in S$、``在 $E$ 中''、$x\in U$、$Y_{\xi,s}$、$F_{\xi,s}$ 等互不协调的写法。按上下文，首处当为 $s\in S$，邻域应在基底 $S$ 中；纤维指标也应与遍历的基底点一致。这里不据此编造一个已经审定的修正版。}

\entry{7.3}{命题}
基本相同。不再有 $F$，并假设 $X$ 相对 $S$ 平坦；以（4.3）的性质(i)至(v)代替（4.2）的那些性质（“但应注意作出保留条件”）：$k(s)$ 的特征为零，或者 $f:X\to P$ 是浸入；而在(v)的情形，对每个 $s\in S$，〔字迹不清〕$X_s$ 的不可约分支 $Z$ 都有
\[
 \dim f(Z)\geq2.
\]
〔附记：对于(v)，比较12.2.1的(x)和(xi)（这时可以〔字迹不清〕在另一种情形4.3或12.2.1的(x)）。〕〔英译编者说明：上述方括号内的页边备注大部分难以辨认。〕

〔英译编者标记：一段文字被划去。〕

\entry{7.4}{命题}
设 $g:X\to S$ 为平坦固有态射，$s\in S$。令 $n=\dim X_s$，并假设 $X_s$ 中 $X_s$ 在 $k(s)$ 上不可分的点所成集合的维数 $\leq n$，例如 $X_s$ 可分。那么，存在 $s$ 的开邻域 $U$，以及平展、有限且满射的态射
\[
 S'\longrightarrow U,
\]
使得 $X\times_S S'$ 在 $S'$ 上有一个截面。如果 $k(s)$ 是无限域，则可以取 $S'$ 为 $X|_U$ 的一个闭子概形。\ednote{不清楚，请问 A.G.。〔底本脚注27。〕}
\begin{proof}
先假设 $k(s)$ 无限。对 $n$ 归纳，$n=0$ 的情形是平凡的。事实上，这时存在 $s$ 的一个开邻域 $U$，使 $X|_U$ 本身在 $U$ 上平展、有限且满射；通过直接引用有关结果即可看出这一点。如果 $n>0$，把（7.3）应用于“可分”这一性质，就能用一个具有同样性质的“超平面截面”$Y$ 来替代 $X$，只是把 $n$ 换成 $n-1$。如果不假设 $k(s)$ 无限，则先作一个平展基变换，论证仍然适用。（可以进行下去。）
\end{proof}
\tnote{底本写的是 $\leq n$，不是 $<n$。这一条件不能支撑其写出的 $n=0$ 起始步。例如 $X=\Spec(k[\varepsilon]/(\varepsilon^2))$、$S=\Spec k$，其中 $k$ 为无限完美域，则 $X\to S$ 平坦且有限，故固有；$n=0$，不可分点集的维数至多为0。但这个非约化的 $X$ 不在 $k$ 上平展。因此“$X|_U$ 本身平展”的推断不成立。此例只反驳所写证明的这一推断，并不反驳这个例子中的截面存在结论。另取无限的不完美域 $k=\mathbf F_p(t)$ 与 $L=k(t^{1/p})$，令 $X=\Spec L$；同样满足所写的维数上界，而不存在非空有限平展 $k$-概形 $S'$ 使 $X_{S'}\to S'$ 有截面，详见本批后附核查说明。故所印条件下的命题本身也不能直接使用。}

\entry{7.5}{备注}
特别地，如果 $X$ 在 $S$ 上射影且可分，则它在 $S$ 上局部地容许平展多截面。不过要注意，可以给出 $X$ 在 $S$ 上固有且光滑、但非射影的例子，使同样结论不成立。当然，一般不能把射影假设减弱为拟射影假设；例如取 $X$ 在 $S$ 上平展但非有限，就可看出这一点……\ednote{字迹不清。〔底本脚注28。〕}
\tnote{（7.4）只明写固有性，而（7.5）又强调非射影情形可以失败；这反映出草稿上下文或假设尚未定型。必须在确定“局部”“多截面”的精确意义及前文嵌入条件后才能整理成定理，不能把这些句子当作已完成的统一证明。}

\backmatter
\chapter*{本批核查说明}
\noindent\textbf{（7.4）按字面的反例。}
取 $k=\mathbf F_p(t)$ 及 $L=k(u)$，其中 $u^p=t$。扩张 $L/k$ 是次数 $p$ 的非平凡纯不可分扩张。$X=\Spec L\to\Spec k$ 有限、平坦、固有，且 $\dim X=0$；其不可分点集也是维数0，所以符合底本的 $\leq n$。

若存在有限平展满射 $S'\to\Spec k$，则 $S'=\Spec A$，其中 $A$ 为非零有限可分 $k$-代数，即有限个有限可分域扩张 $K_i/k$ 的直积。$X_{S'}\to S'$ 的截面会诱导一个 $k$-代数同态 $L\to A$，再投影到任一 $K_i$，得到 $k$-嵌入 $L\hookrightarrow K_i$。这是不可能的，因为有限可分扩张的子扩张可分，而 $L/k$ 非平凡纯不可分。由此，所印条件不能保证结论。本批不自行把 $\leq n$ 改成 $<n$，也不声称仅改变这一符号就能修复全部证明。

\bigskip
\noindent\textbf{覆盖与检查。}
本批逐段核对底本第19--21页的第7节，包括（7.1）两项结论、（7.2）、（7.3）的未辨字及划去标记、（7.4）陈述和证明、（7.5）及脚注26--28；未将第8节计入完成。数学记号与底本图像核对，疑点另列。未核对法文原手稿，未重新审定前文所引用定理；本批数学说明仅针对明确列出的两处反例及文字问题。

\bigskip
\noindent\textbf{下一接续点：}第一篇第8节（8.1）。第三版合订本的其他未完成部分仍保持未完成状态。
\end{document}
